\documentclass[12pt]{article}

\input{../../_assets/preambulo.tex}


\begin{document}

    % 1. Foto de fondo
    % 2. Título
    % 3. Encabezado Izquierdo
    % 4. Color de fondo
    % 5. Coord x del titulo
    % 6. Coord y del titulo
    % 7. Fecha

    
    \input{../../_assets/portada}
    \portadaExamen{etsiitA4.jpg}{LMD\\Examen VI}{LMD. Examen VI}{MidnightBlue}{-8}{28}{2026}{Roxana Acedo Parra}

    \begin{description}
        \item[Asignatura] Lógica y métodos discretos.
        \item[Curso Académico] 2025-26.
        \item[Grado] Doble Grado en Ingeniería Informática y Matemáticas.
        \item[Grupo] Único.
        \item[Profesor] Francisco Miguel García Olmedo.
        \item[Descripción] Final extraordinario.
        \item[Fecha] 17 de julio de 2026.
        \item[Duración] 200 minutos.
    \end{description}
    \newpage
    
    \begin{ejercicio} Basándose en la teoría de recurrencias lineales demuestre que $1=0.\hat{9}$.
    \end{ejercicio}

    
    \begin{ejercicio} Haga razonadamente lo siguiente:
        \begin{enumerate}[label=\alph*)]
            \item Dé dos ejemplos de álgebra de Boole $A$ con 256 elementos.
            \item Demuestre que cualquier álgebra de Boole $A$ con 256 elementos es isomorfa al álgebra de Boole producto $B \times C$, siendo $B$ (resp. $C$) un álgebra de Boole con 3 (resp. 5) átomos.
        \end{enumerate}
    \end{ejercicio}

    
    \begin{ejercicio} Sea $\cc{L}$ un lenguaje de primer orden, $\Gamma$ un conjunto de fórmulas de $\cc{L}$, $c$ un símbolo de constante de $\cc{L}$ que no ocurre en los elementos de $\Gamma$ y $y$ un símbolo de variable de $\cc{L}$. Para toda fórmula $\varphi$, sea $\varphi_c^y$ la fórmula que resulta de sustituir por $c$ cada ocurrencia libre de $y$ en $\varphi$ y sea $\Gamma_c^y = \{\varphi_c^y : \varphi \in \Gamma\}$. Demuestre que son equivalentes las siguientes afirmaciones:
    \begin{enumerate}[label=\alph*)]
        \item $\Gamma$ es insatisfacible.
        \item $\Gamma_c^y$ es insatisfacible.
    \end{enumerate}
    \end{ejercicio}


    \begin{ejercicio} Sabemos que un monoide $M=\langle M, \cdot, e\rangle$ es un álgebra de tipo $\langle 2,0 \rangle$ que cumple las siguientes propiedades:
        \begin{enumerate}[label=\alph*)]
            \item Para todo $x, y, z \in M$, $x \cdot (y \cdot z) = (x \cdot y) \cdot z$ (asociatividad)
            \item Para todo $x \in M$, $x \cdot e = x = e \cdot x$ (propiedad de elemento neutro)
        \end{enumerate}
        Usando las técnicas de resolución lineal input estudiadas en la asignatura, demuestre que todo monoide $M=\langle M, \cdot, e\rangle$ en el que se cumpla la propiedad:
        $$\text{para todo } x \in M, \quad x \cdot x = e$$
        debe ser un monoide conmutativo, i.e., para todo $x, y \in M$ se cumple $x \cdot y = y \cdot x$.
    \end{ejercicio}

    
    \begin{ejercicio} La siguiente matriz, $A$, debe ser la matriz de adyacencia de un grafo con 6 vértices, simple, sin lazos y no dirigido. Complétela razonadamente como sea conveniente en cada caso para que el grafo resultante sea como se pide en cada apartado:
        \begin{enumerate}[label=\alph*)]
            \item un árbol con 3 hojas (vértices de grado 1).
            \item regular de grado 2.
            \item conexo, no sea un árbol y no sea de Euler.
            \item tenga el mayor número posible de lados.
            \item tenga número cromático 3.
        \end{enumerate}
        $$
        A = \begin{pmatrix}
        0 & 1 & 0 & 0 & 0 & * \\
        1 & 0 & 1 & 0 & 0 & * \\
        0 & 1 & 0 & 1 & 0 & * \\
        0 & 0 & 1 & 0 & 1 & * \\
        0 & 0 & 0 & 1 & 0 & * \\
        * & * & * & * & * & *
        \end{pmatrix}
        $$
    \end{ejercicio}

\end{document}